\subsection{Overview}

The implementation is divided into two connected artifacts: a \textbf{research pipeline} that trains and validates flow classifiers on secondary datasets, and a \textbf{Windows desktop app} that deploys the best ECG-only models for live League of Legends sessions. The research pipeline is implemented in Python under \texttt{C:/Users/user/Downloads/Masterthesis/Code}; the app is implemented under \texttt{.../Code/app}. This chapter describes both artifacts, with emphasis on the design decisions that bridge offline benchmarking and live deployment.

\subsection{Research Pipeline: Datasets and Labeling}

\subsubsection{Datasets}

Two secondary datasets are used for model development:

\begin{itemize}
    \item \textbf{BIRAFFE2} \autocite{kutt2022biraffe2}: 102 subjects playing an affective game while ECG, EDA, EEG, and webcam footage are recorded. The flow label is derived from the Game Experience Questionnaire (GEQ) Core Module, specifically the five-item GEQ-R 2018 Flow subscale \autocite{ijsselsteijn2013game}. The 2018 revised scoring is preferred over the 2013 version because it removes cross-loading items and is the current best practice.
    \item \textbf{Irshad/PhySF} \autocite{irshad2023wearable}: 25 subjects performing work-like tasks while ECG, EDA, and EEG are recorded. Flow labels are binary and come from retrospective self-reports about in which third of the task flow occurred.
\end{itemize}

\subsubsection{Labeling Strategies}

Several labeling strategies are compared to understand how much the apparent predictability of flow depends on label construction:

\begin{itemize}
    \item \textbf{Single-level median split:} use only the level-1 GEQ Flow score; classify windows as high or low flow relative to the median.
    \item \textbf{Averaged levels:} average the three GEQ Flow scores across levels and split at the median.
    \item \textbf{Levels as pseudo-subjects:} split the continuous GAME-phase recording into three equal segments, assign each segment its corresponding level score, and treat each segment as a separate pseudo-subject in leave-one-subject-out (LOSO) cross-validation. This design approximates time-varying labels without continuous self-report interruptions.
    \item \textbf{Raw-item recomputation:} recompute the Flow score from the raw GEQ item responses, optionally dropping the Time Distortion item (item 25) to test whether retrospective time perception adds signal.
    \item \textbf{Extreme percentile:} label only the bottom and top 20\% of scores as low/high flow, removing ambiguous near-median subjects.
\end{itemize}

\subsubsection{Preprocessing and Windowing}

ECG is cleaned with NeuroKit2 \autocite{makowski2021neurokit2} and converted to R-R intervals. From these intervals, HRV features are computed in 60-second sliding windows with a 30-second step (50\% overlap). This windowing matches the best offline setup and is reproduced in the live app. EDA and EEG features are extracted for offline multimodal benchmarks but are not deployed live. Z-standardisation is applied per cross-validation fold, and an optional IQR outlier handler is evaluated.

\subsection{Research Pipeline: Feature Extraction}

For each 60-second ECG window the following features are extracted:

\begin{itemize}
    \item \textbf{Time domain:} mean heart rate (\texttt{hr\_mean}), SDNN, RMSSD, pNN50.
    \item \textbf{Frequency domain:} VLF, LF, HF, total power (TP), and LF/HF ratio.
    \item \textbf{Non-linear:} sample entropy, DFA $\alpha_1$, DFA $\alpha_2$.
\end{itemize}

Because VLF and DFA $\alpha_2$ require several minutes of data, they are often NaN for 60-second windows. A diagnostic setup (01d) drops these unreliable features to test whether the shorter-window feature set is cleaner; it is not. The final app feature extractor computes the same 12-feature vector and fills any remaining all-NaN columns with zeros.

\subsection{Research Pipeline: Models and Validation}

\subsubsection{Classifiers}

A multi-model benchmark compares classical and deep-learning classifiers:

\begin{itemize}
    \item \textbf{Classical:} Random Forest, SVM (RBF), k-Nearest Neighbors, Logistic Regression, XGBoost.
    \item \textbf{Deep learning:} MLP, LSTM, 1D-CNN (all implemented in PyTorch and consuming the same flat per-window feature matrix).
\end{itemize}

Classical hyperparameters are documented in \texttt{results/hyperparameters.md}: Random Forest uses 200 estimators with balanced class weights; SVM uses an RBF kernel with probability estimates; kNN uses $k=5$; XGBoost uses 200 estimators, max depth 4, and learning rate 0.05.

\subsubsection{Validation Strategy}

Models are evaluated with \textit{Leave-One-Subject-Out} (LOSO) cross-validation. To avoid data leakage, overlapping windows from the same subject never appear in both training and test sets. Because some folds have only one class in the training set, a subject-level aggregation is added: all windows of a left-out subject are classified, and the majority prediction determines the subject-level label. This allows AUC computation even when individual windows are class-skewed.

\subsection{Research Pipeline: Key Results}

Table~\ref{tab:main_results} summarizes the most important offline results. All metrics are subject-level aggregates from LOSO cross-validation.

\begin{table}[htbp]
    \centering
    \scriptsize
    \caption{Summary of the most important offline classification results.}
    \renewcommand{\arraystretch}{1.3}
    \begin{tabularx}{\linewidth}{|l|l|L|l|l|l|l|}
        \hline
        \rowcolor{black}
        \textbf{\color{white} Setup} & \textbf{\color{white} Modality} & \textbf{\color{white} Best Model} & \textbf{\color{white} Acc} & \textbf{\color{white} F1} & \textbf{\color{white} AUC} & \textbf{\color{white} n} \\
        \hline
        01 single-level & ECG & CNN1D & 0.394 & 0.383 & 0.290 & 99 \\
        01c levels-as-subjects & ECG & CNN1D & 0.538 & 0.528 & 0.556 & 247 \\
        06 multimodal & ECG+EDA+FACE & CNN1D & 0.564 & 0.554 & 0.561 & 273 \\
        06 multimodal (grid) & ECG+EDA+FACE & RandomForest & 0.612 & 0.598 & 0.623 & 268 \\
        09 heartpy & ECG & MLP & 0.420 & 0.408 & 0.376 & 100 \\
        12 baseline correction & ECG & MLP & 0.628 & 0.627 & 0.677 & 223 \\
        12 baseline correction (grid) & ECG & RandomForest & 0.647 & 0.647 & 0.681 & 207 \\
        12 extreme percentile & ECG & SVM & 0.730 & 0.727 & 0.790 & 111 \\
        08b Irshad ECG-only & ECG & RandomForest & 0.680 & 0.667 & 0.712 & 25 \\
        \hline
    \end{tabularx}
    \label{tab:main_results}
\end{table}

The main findings are:

\begin{itemize}
    \item \textbf{Pure ECG performs poorly with naive labels.} Setup 01 (single-level median split) reaches only AUC = 0.290, indicating that subject-specific baseline physiology dominates between-subject HRV differences.
    \item \textbf{Levels-as-subjects improves signal.} Treating each game level as a pseudo-subject raises AUC to 0.556, because within-subject variation is more informative than between-subject averages.
    \item \textbf{Baseline correction is the strongest single improvement.} Setup 12 subtracts a resting baseline from each feature and reaches AUC = 0.677, suggesting that removing subject-specific resting physiology unmasks a within-session flow signal.
    \item \textbf{Multimodality helps only marginally.} Adding EDA and webcam affect (Setup 06) reaches AUC = 0.561 without grid search and 0.623 with grid search, but the improvement is smaller than the baseline-correction effect and the setup is more complex to deploy.
    \item \textbf{Irshad/PhySF ECG-only is competitive.} The Irshad RandomForest ECG-only model reaches AUC = 0.712 on 25 subjects, providing a cross-dataset validation of the ECG-only approach.
\end{itemize}

Permutation importance analysis confirms that the most informative ECG features are RMSSD, SDNN, LF, and pNN50, with no single feature dominating. Feature-group permutation shows that ECG contributes more than EDA or webcam affect in the multimodal setup, reinforcing the decision to deploy ECG-only.

\subsection{Live App: Architecture}

The live app is a PyQt6 Windows desktop application packaged with PyInstaller and Inno Setup. It follows a local-first architecture:

\begin{itemize}
    \item \textbf{Frontend:} PyQt6 main window with sidebar navigation (Dashboard, Logs/Webcam, Settings), system tray icon, and toggle-button controls.
    \item \textbf{Sensors:} Polar H10 BLE chest belt via \texttt{polar\_python}; optional webcam background recording via OpenCV.
    \item \textbf{Inference:} pre-trained ECG-only classifiers loaded from disk; hard majority vote.
    \item \textbf{Game integration:} local League Client (LCU) API for match start/end detection, queue type, and live event polling.
    \item \textbf{Persistence:} local SQLite database and gzipped CSV files for ECG; MP4 files for webcam.
\end{itemize}

\subsection{Live App: Sensor Layer}

\subsubsection{ECG Streaming}

The Polar H10 chest belt streams ECG over Bluetooth Low Energy using the Polar PMD protocol. The app uses \texttt{polar\_python} (built on \texttt{bleak}) to handle the control point and ECG notifications. The implementation supports:

\begin{itemize}
    \item reactive pairing if Windows reports insufficient authentication;
    \item direct connection to a saved MAC address, bypassing scanner issues on Windows;
    \item buffering of incoming samples until a 60-second window is complete.
\end{itemize}

A 30-second overlap buffer is maintained so a new prediction is produced every 30 seconds, matching the offline pipeline.

\subsubsection{Webcam Recording}

The webcam records an MP4 file per session using OpenCV with an \texttt{avc1} codec and an \texttt{mpv4} fallback. Recording starts only after a sensor connects, and the user can enable or disable it at any time. The video is
\textit{not} fed into the live classifier; it is stored for post-hoc reflection and qualitative validation.

\subsection{Live App: Inference Layer}

\subsubsection{Exported Model Ensemble}

The app loads five ECG-only models exported from the research pipeline:

\begin{table}[htbp]
    \centering
    \scriptsize
    \caption{Models exported to the live app ensemble.}
    \renewcommand{\arraystretch}{1.3}
    \begin{tabularx}{\linewidth}{|l|l|l|l|l|}
        \hline
        \rowcolor{black}
        \textbf{\color{white} Model File} & \textbf{\color{white} Dataset} & \textbf{\color{white} Type} & \textbf{\color{white} AUC} & \textbf{\color{white} Notes} \\
        \hline
        \texttt{biraffe2\_svm.joblib} & BIRAFFE2 baseline-corrected & classical & 0.790 & Best extreme-percentile model \\
        \texttt{biraffe2\_knn.joblib} & BIRAFFE2 baseline-corrected & classical & 0.761 & Strong non-linear baseline \\
        \texttt{biraffe2\_rf.joblib} & BIRAFFE2 baseline-corrected & classical & 0.725 & Robust ensemble member \\
        \texttt{irshad\_rf.joblib} & Irshad/PhySF ECG-only & classical & 0.692 & Cross-dataset validation \\
        \texttt{biraffe2\_mlp.joblib} & BIRAFFE2 baseline-corrected & deep & 0.730 & Exported as weight/scaler dict \\
        \hline
    \end{tabularx}
    \label{tab:app_models}
\end{table}

The AUC values in Table~\ref{tab:app_models} come from the best offline configurations used for export and represent the intended ensemble composition. During live use, each model classifies the current 60-second window, and the ensemble label is the hard majority vote. With five binary classifiers, a majority is always guaranteed when an odd number of models is enabled.

\subsubsection{Feature Extraction at Runtime}

The runtime feature extractor reimplements the offline HRV extraction path using the same NeuroKit2 functions. For each completed 60-second window it computes the 12 ECG features, imputes all-NaN columns with zero, applies the scaler saved during model export, and passes the vector to each classifier.

\subsection{Live App: Game Integration}

Match detection follows a fallback chain:

\begin{enumerate}
    \item \textbf{LCU API first:} read the local League Client lockfile, authenticate to the local LCU API, and poll gameflow state to detect match start/end and queue type (Solo/Duo, Flex, ARAM, etc.).
    \item \textbf{Process detection:} if the LCU API is unavailable, detect the running \texttt{League of Legends.exe} process.
    \item \textbf{Manual control:} the user can start or stop a match manually at any time.
\end{enumerate}

During a match, the app polls the \textbf{Riot Games Live Client Data API} (\url{https://127.0.0.1:2999/liveclientdata/eventdata}) every 1--2 seconds to capture events such as ChampionKill, TurretKilled, DragonKilled, and BaronKilled. These events are stored in the local SQLite database and later displayed on the post-game timeline.

\subsection{Live App: Adaptation Layer}

The current adaptation layer is entirely post-game:

\begin{itemize}
    \item \textbf{Flow timeline:} a per-window flow/non-flow trajectory is drawn together with event markers (green for self kills, red for self deaths, blue for objectives).
    \item \textbf{Match summary:} a rule-based analyzer computes the percentage of flow windows, the early-game versus late-game trend, and the correlation between deaths and no-flow windows. It outputs one of four recommendations: Continue, ShortBreak, LongBreak, or WarmUp.
    \item \textbf{Replay linking:} the app scans \texttt{\%USERPROFILE\%/Documents/League of Legends/Replays}, reads the \texttt{.rofl} metadata header, and links replays to recorded sessions by timestamp. A reflection prompt can direct the player to the replay segment around a detected no-flow/death event.
\end{itemize}

\subsection{Live App: Persistence and Privacy}

All data remain on the local machine. The SQLite database stores sessions, matches, predictions, game events, and sensor events. ECG is written as gzipped CSV, the webcam as MP4, and logs as text files. The user can choose storage paths in the Settings dialog and delete individual sessions from the dashboard.

\subsection{Deployment}

The app is built with:

\begin{itemize}
    \item \texttt{scripts/export\_models.py} trains and exports the five model bundles from the research pipeline.
    \item \texttt{scripts/build\_app.py} runs PyInstaller to produce a single-folder or one-file distribution.
    \item \texttt{scripts/installer.iss} is an Inno Setup script that creates \texttt{FlowLoL\_Setup.exe} for lab PCs.
\end{itemize}

The installed app stores user data under \texttt{\%LOCALAPPDATA\%/FlowLoL/Data} with user-modify permissions, so ECG and webcam recordings can be written without administrator rights.

\subsection{Summary}

The implementation connects an offline research pipeline with a live Windows companion app. The pipeline establishes that ECG-only models with baseline correction provide the best deployment-ready signal, while the app turns that signal into a privacy-preserving, post-game reflection tool. The most ambitious real-time interventions identified in the qualitative requirements are deliberately deferred until the live classification signal is validated with real players.
